\subsection*{定义有界变量}

\co@with-in@ 可以定义有界整数或有界整数元组，不创建额外的执行实例：

\begin{lstlisting}[language=croqtile,caption={一个有界变量}]
with x in 128 {
  // 'x' is a bounded integer
}
with y in [512] {
  // 'y' is a bounded ituple
}
with index in [10, 10] {
  // 'index' is a bounded ituple
}
\end{lstlisting}

本例中，\co@x@ 的范围是 128；\co@y@ 是一维元组，范围为 \co@[512]@；\co@index@ 是二维元组，范围为 \co@[10, 10]@。元组的元素也可以分别命名。

\begin{lstlisting}[language=croqtile,caption={一组有界变量}]
with {x, y} in [10, 10] {
  // x and y are explicitly named elements of the ituple
}
with index = {x, y} in [10, 10] {
  // 'x', 'y', and 'index' can be used within the block
}
\end{lstlisting}

多条声明可以用逗号合写：

\begin{lstlisting}[language=croqtile,caption={由元组导出的变量}]
with index = {x, y} in [10, 10], idx in [100, 10] { }
\end{lstlisting}

\co@index@ 与 \co@idx@ 只在随后的 \co@with-in@ 作用域内有效。
